\subsection{A five-object model and the exact object minimum}


\begin{remark*}[A singleton extreme hom-set and the fundamental group]
Suppose a category has objects $s,t$ such that every object $x$ admits
arrows $a_x:s\to x$ and $b_x:x\to t$, while $C(s,t)=\{e\}$.
Then its classifying space is simply connected. Indeed, choose
$a_s=1_s$ and $a_t=e$. In the groupoid obtained by inverting every
arrow, an arrow $f:x\to y$ satisfies
\[
 b_y f a_x=e=b_y a_y,
 \qquad\text{hence}\qquad f=a_y a_x^{-1}.
\]
Thus every edge path from $s$ to $x$ equals $a_x$, and every loop
at $s$ is trivial. This groupoid is the fundamental groupoid of
the nerve: its presentation has the category arrows as edges and
one relation $[g][f]=[g\circ f]$ for each two-simplex. All four
models below have a singleton $C(0,4)$ and full linear support,
so this gives a direct proof of their simple connectivity in
addition to the independently checked Tietze transcripts.
\end{remark*}

\begin{theorem}[Proved]\label{A7:model47}
There is a finite acyclic category with five objects and forty-seven
arrows, identities included, whose classifying space is homotopy
equivalent to $\mathbb{CP}^2$.
\end{theorem}
\begin{proof}
We give the entire category. Let $M$ be the first explicit monoid
of Theorem~\ref{A:monoid-upper23}, with its numeric element labels
$0=1$, $x_{ij}=1+5i+j$, and actors $16,\ldots,22$.
Our objects are $0<1<2<3<4$. Arrows $0,\ldots,4$ are their
identities. The nonidentity arrows and an auxiliary monoid label
$\ell$ are as follows; lists correspond in the displayed order.
\[
\begin{array}{c|c|l}
\text{hom-set}&\text{arrow labels}&\text{monoid labels }\ell\\\hline
01&5,6&1,3\\
02&7,8,9,10&1,2,5,3\\
03&11,12,13&1,4,5\\
04&14&1\\
12&15,\ldots,21&1,2,11,16,17,20,21\\
13&22,\ldots,34&1,1,1,4,4,5,11,11,11,14,14,19,21\\
14&35,36&1,11\\
23&37,\ldots,41&1,4,18,19,22\\
24&42,43,44&1,6,11\\
34&45,46&1,6
\end{array}
\]
Identities have their usual products and auxiliary label $0$.
For $b:1\to2$ and $c:2\to3$, the product $c\circ b$ is the
following entry, with rows $b$ and columns $c$:
\[
\begin{array}{c|rrrrr}
 &37&38&39&40&41\\\hline
15&22&25&23&22&24\\
16&24&26&27&25&22\\
17&28&31&29&30&29\\
18&24&26&24&33&33\\
19&23&25&27&33&34\\
20&28&31&34&31&33\\
21&29&32&34&29&34
\end{array}
\]
Every other composable pair of nonidentity arrows
$a:i\to j$, $b:j\to k$ has as product the unique arrow
in $C(i,k)$ with auxiliary label $\ell(a)\ell(b)$, multiplied
in $M$ in this chronological order. All such target hom-sets
have distinct auxiliary labels. Inspection of the monoid table
checks existence of the indicated target and that every entry
of the exceptional table has the prescribed product label.
These rules specify every composable pair and no untyped pair.

Any triple of nonidentity arrows has total object-distance at
least three. Its target hom-set has distinct auxiliary labels,
so the two bracketings agree by associativity in $M$. Thus these
rules define a category; its support and identity endomorphisms
make it acyclic. Its hom profile is
$(2,4,3,1,7,13,2,5,3,2)$, of sum $42$, hence the arrow count
is $47$. A literal $47\times47$ partial composition table is in
\path{acyclic/outputs/cp2_five_47_integral_certificate.json}.

Here are explicit cohomological data. On chronological
two-simplices $(a,b)$ define the integral cochain $\alpha$ to be
zero except in the following cases:
\[
\begin{array}{c|l|r}
a&b&\alpha(a,b)\\\hline
6&17,20,21,28,29,30,31,32,33,34,36&1\\
8,16,20&40,41,43,44&1\\
9,19,21&43&1\\
10&40,41,44&1\\
13,27,34&46&1\\
12,25,26,31,32,38&46&-1
\end{array}
\]
Every pair specified in this table is composable. Let
\[
 y=[5,15]-[5,17]-[6,15]+[6,17].
\]
Let $W$ be the integer matrix below, whose rows are indexed by
$15,\ldots,21$ and columns by $37,\ldots,41$:
\[
W=\begin{pmatrix}
-2&1&1&1&-1\\
1&0&-1&-1&1\\
1&-1&0&0&0\\
1&0&-1&-1&1\\
-1&0&1&1&-1\\
-1&0&1&1&-1\\
1&0&-1&-1&1
\end{pmatrix}.
\]
Define the four-chain
\[
 z=\sum_{b=15}^{21}\sum_{c=37}^{41}W_{bc}
 \bigl([5,b,c,45]-[5,b,c,46]
       -[6,b,c,45]+[6,b,c,46]\bigr).
\]
The indices on $W$ here mean the corresponding row and column.
Direct substitution in the five signed nerve faces verifies
\[
 \delta\alpha=0,\qquad \partial y=0,\qquad \partial z=0,
 \qquad \langle\alpha,y\rangle
 =\langle\alpha\smile\alpha,z\rangle=1.
\]
The program \path{acyclic/code/check_cp2_five_integral.py}
imports no constructor or search code, reconstructs the full
normalized nerve, verifies every displayed equation and all
$927$ associativity equations, and checks the $274$ cocycle
equations integrally.

The integral nerve has cell counts
\[
 (f_0,f_1,f_2,f_3,f_4)=(5,42,174,274,140).
\]
The four boundary ranks are $(4,38,135,139)$ and every nonzero
Smith invariant is $1$. These statements are checked by the
preceding program and independently by
\path{monoid/code/check_acyclic47.py}, with no acyclic
producer imported. The full latter certificate is recorded with
\path{monoid/certificates/acyclic47_integral_independent_verified.json}.
For clarity, Smith invariants of the full boundaries suffice:
$\ker\partial_d$ is saturated in the free group $C_d$, so
$C_d/\operatorname{im}\partial_{d+1}$ is an extension of $H_d$
by the free group $\operatorname{im}\partial_d$. The extension
splits, and the torsion of $H_d$ is exactly the torsion of this
boundary cokernel. The ranks therefore give
\[
 H_*(BC;\mathbb Z)=(\mathbb Z,0,\mathbb Z,0,\mathbb Z),
\]
with zero groups above degree four by acyclicity. The displayed
pairings make $[\alpha]$ a generator and its square a generator.

For simple connectivity, collapse the spanning tree formed by
arrows $5,15,37,45$. The edge-and-triangle presentation has
$38$ generators and $174$ relations. A transcript of $38$
single-occurrence Tietze eliminations reduces it to the empty
presentation. The independent
\path{code/check_acyclic_fundamental_group.py} reconstructs the
typed table and presentation and checks every substitution in
\path{outputs/five_object_47_pi1.json}; its verification output
is \path{outputs/five_object_47_pi1.json.verified.json}. Consequently
$\pi_1(BC)=1$.

We can now apply the homotopy classification in
\cite[Proposition~4C.3, pp.~430--431]{CP2-Hatcher2002}, whose
statement and proof were opened in this run. Its hypotheses are
simple connectivity, CW type, and reduced integral homology
$H_2=\mathbb Z^m$, $H_4=\mathbb Z$, with all other groups zero;
it says the integral cohomology ring determines the homotopy
type. Here $m=1$ and the primitive square gives precisely the
ring of $\mathbb{CP}^2$. Thus $BC\simeq\mathbb{CP}^2$.
\end{proof}

\begin{corollary}[Proved]\label{A7:nu}
The least number of objects of a finite acyclic category of type
$\mathbb{CP}^2$ is exactly $\nu(\mathbb{CP}^2)=5$.
\end{corollary}
\begin{proof}
Proposition~\ref{A7:dimension-chain} gives the lower bound, and
Theorem~\ref{A7:model47} gives a model with five objects.
\end{proof}

\begin{theorem}[Proved]\label{A7:model42}
There is a five-object acyclic model of $\mathbb{CP}^2$ with
forty-two arrows, identities included.
\end{theorem}
\begin{proof}
Use the same monoid and composition convention as in
Theorem~\ref{A7:model47}. Objects and identities are $0,\ldots,4$.
The complete nonidentity hom-sets and auxiliary monoid labels are
\[
\begin{array}{c|c|l}
\text{hom-set}&\text{arrow labels}&\ell\\\hline
01&5,6&1,3\\
02&7,8,9,10&1,2,5,3\\
03&11,12,13&1,4,5\\
04&14&1\\
12&15,\ldots,21&1,2,11,16,17,20,21\\
13&22,\ldots,30&4,4,1,5,14,14,11,19,21\\
14&31,32&1,11\\
23&33,34,35,36&4,18,19,22\\
24&37,38,39&1,6,11\\
34&40,41&1,6
\end{array}
\]
The exceptional composition table, again for $1\to2\to3$, is
\[
\begin{array}{c|rrrr}
&33&34&35&36\\\hline
15&23&24&24&24\\
16&22&25&22&24\\
17&26&28&28&28\\
18&23&24&29&29\\
19&22&25&29&30\\
20&26&30&26&29\\
21&27&30&28&30
\end{array}
\]
Every other composite is the unique arrow in its target hom-set
whose auxiliary label is the chronological monoid product.
As before this completely specifies a category and proves
associativity; the independent programs also check all $773$
associativity equations. Its hom profile is
$(2,4,3,1,7,9,2,4,3,2)$.

An integral two-cocycle is zero except in these cases:
\[
\begin{array}{c|l|r}
a&b&\alpha(a,b)\\\hline
6&17,20,21,26,27,28,29,30,32&1\\
8,16,20&35,36,38,39&1\\
9,19,21&38&1\\
10&35,36,39&1\\
13,25,30&41&1\\
12,22,23,26,27,33&41&-1
\end{array}
\]
Set $y=[5,16]-[5,17]-[6,16]+[6,17]$, and let
\[
W=\begin{pmatrix}
0&0&0&0\\
1&-1&-1&1\\
-1&1&1&-1\\
0&-1&-1&1\\
0&1&1&-1\\
0&1&1&-1\\
0&-1&-1&1
\end{pmatrix},
\]
with rows $15,\ldots,21$ and columns $33,\ldots,36$. Define
\[
 z=\sum_{b,c}W_{bc}
  \bigl([5,b,c,40]-[5,b,c,41]
        -[6,b,c,40]+[6,b,c,41]\bigr).
\]
The literal integral checks give $\delta\alpha=0$,
$\partial y=\partial z=0$, and
$\langle\alpha,y\rangle=\langle\alpha\smile\alpha,z\rangle=1$.
The normalized nerve has counts $(5,37,145,222,112)$ and boundary
ranks $(4,33,111,111)$; every nonzero Smith invariant is $1$.
Thus its integral homology is $(\mathbb Z,0,\mathbb Z,0,\mathbb Z)$,
with no higher cells. These facts are checked independently in
\path{acyclic/outputs/cp2_five_42_integral_certificate.json} and
\path{monoid/certificates/acyclic42_integral_independent_verified.json}.
Both checks reconstruct every chain and use exact integer matrices.

The spanning tree $5,15,33,40$ leaves $33$ generators and $145$
triangle relations. The same independent fundamental-group
checker verifies $33$ elementary Tietze eliminations to the empty
presentation, with transcript
\path{acyclic/outputs/cp2_five_42_pi1.json} and result
\path{acyclic/outputs/cp2_five_42_pi1.json.verified.json}.
Hence the space is simply connected. Proposition~4C.3 of
\cite{CP2-Hatcher2002} now applies exactly as in
Theorem~\ref{A7:model47}, proving its type.
The complete partial table and all literal cochains and cycles
are in the stated integral certificate.
\end{proof}

\begin{remark*}[The sink, join, and product tools]
The forty-two-arrow category is obtained from its full subcategory
$D$ on $0123$ by adding the sink $4$ along the presheaf
$F:D^{\mathrm{op}}\to\mathrm{FinSet}$, $F(i)=C(i,4)$,
where $F(f)(a)=a\circ f$. Its stalk sizes are $(1,2,3,2)$.
Thus the sink lemma, \cite[Lemma~2.1]{CP2-Carcacia2026}, identifies
its classifying space with the homotopy cofiber of the explicitly
specified projection $B\operatorname{El}(F)\to BD$. This is a
second, attachment-based description of the same fully given
category; the integral checks above compute the resulting type.

The join lemma, \cite[Lemma~1.2]{CP2-Carcacia2026}, also provides
a useful obstruction: no categorical join of two nonempty finite
categories can model $\mathbb{CP}^2$. Indeed $B(C\star D)$ is
$BC*BD$. After choosing basepoints, collapse the contractible
subcomplex consisting of the two basepoint cones. The quotient
is the reduced suspension of the smash product. Its positive
cohomology products vanish, whereas the degree-two square on
$\mathbb{CP}^2$ does not. This explains why the five-object
join constructions that suffice for spheres cannot give this
model. Finally, the same lemma gives
$B(C\times D)=BC\times BD$ and the product object bound.
Applied to two three-object sphere categories it gives a
nine-object model of $S^2\times S^2$, whose second homology
has rank two. Passing to a categorical symmetric quotient
requires a separate comparison; the explicit quotient checks
reported above show why the product bound alone does not
supply a projective-plane model.
\end{remark*}

\begin{theorem}[Proved: a thirty-nine-arrow model]
\label{A7:model39}
There is a finite acyclic category with five objects and thirty-nine
arrows whose classifying space is homotopy equivalent to
$\mathbb{CP}^2$.
\end{theorem}
\begin{proof}
Let $M$ be the explicit monoid of Theorem~\ref{A:monoid-upper23},
with elements numbered $0,\ldots,22$ and identity $0$.
Take objects $0<1<2<3<4$ and the following nonidentity arrows.
In each row, the arrow numbers and monoid labels correspond in order.
\[
\begin{array}{c|c|l}
(i,j)&\text{arrow numbers}&\text{monoid labels}\\\hline
(0,1)&5,6&1,3\\
(0,2)&7,8,9,10&1,2,3,5\\
(0,3)&11,12,13&1,4,5\\
(0,4)&14&1\\
(1,2)&15,16,17,18,19,20&1,11,16,17,20,21\\
(1,3)&21,22,23,24,25,26,27&1,5,11,14,15,19,21\\
(1,4)&28,29&1,11\\
(2,3)&30,31,32,33&18,19,21,22\\
(2,4)&34,35,36&1,6,11\\
(3,4)&37,38&1,6
\end{array}
\]
Add the five identity arrows, numbered by their objects and labelled
by $0$. There are no other arrows. To compose an arrow with monoid
label $x$ followed by one with label $y$, take the unique arrow in
the required hom-set with monoid label $xy$. The displayed hom-sets
are closed under these products: they are exactly the iterated
products of the four adjacent label sets. Thus this rule specifies
all products, and associativity follows from that of $M$.
The literal checker also verifies every one of the $683$ typed
associativity equations, including identities.

The normalized nerve has simplex counts $(5,34,128,192,96)$ and no
simplices in higher degrees. Its integral boundary ranks are
$(4,30,97,95)$ and every nonzero Smith invariant is one. Hence
its integral homology is $\mathbb Z$ in degrees $0,2,4$ and zero
otherwise. Here is explicit primitive cup-product data. Define
an integral degree-two cochain $\alpha$ by the following nonzero
values, with all unspecified values zero:
\[
\begin{array}{c|l|r}
a&b&\alpha([a,b])\\\hline
6&16,19,20,23,24,25,26,27,29&1\\
8,19&31,32,33,35,36&1\\
9&31,32,33,36&1\\
10,18,20&35&1\\
12,24&38&-1\\
13,22,25,27,32&38&1
\end{array}
\]
This is the pullback of the verified monoid cocycle and satisfies
all $192$ integral cocycle equations. The integral two-cycle
\[
 y=[5,15]-[5,16]-[6,15]+[6,16]
\]
satisfies $\alpha(y)=1$. Let $W$ be the matrix whose rows are
indexed by arrows $15,\ldots,20$ and columns by $30,\ldots,33$:
\[
 W=\begin{pmatrix}
 0&1&-1&0\\
 0&-1&1&0\\
 -1&0&-1&1\\
 1&0&1&-1\\
 1&0&1&-1\\
 -1&0&-1&1
 \end{pmatrix}.
\]
In the tensor notation for four-simplices, put
\[
 z=(5-6)\otimes W\otimes(37-38).
\]
Direct integral boundary and Alexander--Whitney calculations give
$\partial z=0$ and $\langle\alpha\smile\alpha,z\rangle=1$.
Thus $\alpha$ and its square are primitive in degrees two and four.
These assertions, including all integer matrices and chain vectors,
are certified in
\path{acyclic/outputs/cp2_five_39_integral_certificate.json}.
A second implementation, importing no construction or checker code
from this directory, independently computes saturated integral
kernels and the primitive pairing; its receipt is
\path{monoid/certificates/acyclic39_integral_independent_verified.json}.

For the fundamental group, take the spanning tree with arrows
$5,15,30,37$. The resulting presentation has thirty generators
and $128$ relations. The full transcript contains thirty
single-occurrence Tietze eliminations and ends with the empty
presentation. The independently reconstructed and checked transcript
is \path{acyclic/outputs/cp2_five_39_pi1.json}, with a
\texttt{PASS} receipt in the file of the same name followed by
\texttt{.verified.json}. Hence the nerve is simply connected.
Hatcher's Proposition~4C.3, in the form already checked above,
now gives the asserted homotopy equivalence.
\end{proof}

\begin{theorem}[Proved: a thirty-eight-arrow model]
\label{A7:model38}
There is a finite acyclic category with five objects and thirty-eight
arrows whose classifying space is homotopy equivalent to
$\mathbb{CP}^2$.
\end{theorem}
\begin{proof}
Use the same explicitly numbered monoid $M$ and five ordered objects
as above. The complete list of nonidentity arrows is
\[
\begin{array}{c|c|l}
(i,j)&\text{arrow numbers}&\text{monoid labels}\\\hline
(0,1)&5,6&3,5\\
(0,2)&7,8,9,10&1,2,3,5\\
(0,3)&11,12,13&1,4,5\\
(0,4)&14&1\\
(1,2)&15,16,17,18,19&5,16,17,18,20\\
(1,3)&20,21,22,23,24,25&1,5,11,18,19,21\\
(1,4)&26,27,28&1,6,11\\
(2,3)&29,30,31,32&1,18,21,22\\
(2,4)&33,34,35&1,6,11\\
(3,4)&36,37&6,19
\end{array}
\]
Add identities $0,\ldots,4$, all with monoid label $0$.
Composition is again the unique arrow in the required hom-set
labelled by the product of the two labels in chronological order.
Every longer hom-set is exactly the iterated product set of the
four adjacent ones. Closure and associativity therefore follow
directly from $M$; the separate literal checker also verifies all
$621$ typed associativity equations.

The normalized nerve has simplex counts $(5,33,117,166,80)$.
The integral boundary ranks are $(4,29,87,79)$ and every nonzero
Smith invariant is one. Thus the integral homology is
$\mathbb Z$ in degrees $0,2,4$ and zero in all other degrees.
Define the integral two-cochain $\alpha$ by these nonzero values:
\[
\begin{array}{c|l|r}
a&b&\alpha([a,b])\\\hline
5&19,22,24,25,28&1\\
6&27&1\\
8,19&31,32,34,35&1\\
9&31,32,35&1\\
10,15,17&34&1\\
12&36&-1\\
13,21,25,31&36&1
\end{array}
\]
All omitted values are zero. This is the pulled-back monoid
cocycle, and all $166$ integral cocycle equations hold.
In the formal tensor notation for composable arrow strings, put
\[
\begin{aligned}
y&=(5-6)\otimes(-16+17-18+19),\\
w&=(-29+30+31-32)\otimes(36-37),\\
z&=(5-6)\otimes(-16+17-18+19)
   \otimes(-29+30+31-32)\otimes(36-37).
\end{aligned}
\]
The first two expressions are integral two-cycles; the last is an
integral four-cycle with $64$ nonzero coefficients. Direct boundary
calculations and the Alexander--Whitney formula give
\[
 \alpha(y)=\alpha(w)=1,
 \qquad \langle\alpha\smile\alpha,z\rangle=1.
\]
This proves that the degree-two class and its square are primitive.
The literal category, every nonzero Smith invariant, and the
explicit cochain $\alpha$ and cycle $z$ are recorded in
\path{acyclic/outputs/cp2_five_38_integral_certificate.json}.
The second independent implementation gives the same integral
homology and primitive pairing in
\path{monoid/certificates/acyclic38_integral_independent_verified.json}.

The unique arrow $0\to4$ gives simple connectivity by the direct
groupoid argument at the start of this subsection. Independently,
the spanning tree $5,15,29,36$ gives $29$ generators and $117$
triangle relations. The transcript in
\path{outputs/five_object_38_pi1.json}, with its separate verified
receipt, consists of $29$ elementary Tietze eliminations to the
empty presentation. The integral homology and primitive square
now give the claimed homotopy type by
\cite[Proposition~4C.3]{CP2-Hatcher2002}.
\end{proof}

\begin{remark*}[The subsequent exact minimum]
Let $a_5(\mathbb{CP}^2)$ denote the least arrow count of a
five-object acyclic model. The constructions above first established
upper bounds. The separately completed through-thirty-one exclusion
and the next explicit construction settle the exact value
$a_5(\mathbb{CP}^2)=32$.
\end{remark*}
